%% Vintage-aware reconciliation of official SDG indicators.
%% Elsevier CAS double-column class (cas-dc), compiled with pdflatex + bibtex.
%%
%% Working draft written at H10 track C, 12 September 2026. Sections 1-4 report
%% work that is already measured; the injected-revision sweep that answers RQ3
%% is frozen but not yet executed, and Section 5 says so explicitly.
%%
%% EVERY NUMBER in this file comes from papers/vintage_reconciliation/data/
%% angka-kunci.csv through its metric_id. Do not retype numbers by hand.

\documentclass[a4paper,fleqn]{cas-dc}

\usepackage[authoryear,longnamesfirst]{natbib}

\begin{document}
\let\WriteBookmarks\relax
\def\floatpagepagefraction{1}
\def\textpagefraction{.001}

\shorttitle{Vintage-aware reconciliation of official SDG indicators}
\shortauthors{A. Satria et al.}

\title[mode = title]{Vintage-aware reconciliation of official SDG indicators:
measuring the cost of reproducible recomputation on an open table format}

\author[1]{Ardika Satria}
\cormark[1]
\ead{ardika.satria@sd.itera.ac.id}
\credit{Conceptualization, Methodology, Software, Investigation, Writing -- original draft}

\affiliation[1]{organization={Institut Teknologi Sumatera},
            addressline={Jl. Terusan Ryacudu, Way Huwi},
            city={Lampung Selatan},
            postcode={35365},
            state={Lampung},
            country={Indonesia}}

%% TODO: add the five co-authors, their affiliations and CRediT roles before
%% submission. The author block below is a placeholder and must not be left in.
\author[1]{[Co-author names pending]}
\credit{Investigation, Data curation, Validation}

\cortext[1]{Corresponding author}

\begin{abstract}
Official SDG indicators for Indonesia are published through several legitimate
channels of the national statistical office, and those channels disagree. The
disagreement is not an error: it follows from how official statistics are
produced, through release vintages, methodological change and differences in
granularity. Current practice overwrites the older figure with the newer one, so
a number that was once published can no longer be reproduced, even though
revision guidelines for official statistics require exactly that capability.
This paper treats the indicator vintage as a first-class dimension of an open
table format and measures what keeping it costs. Four storage treatments are
compared on an identical revision workload: overwrite, full snapshot, single
source selection under a frozen trust score, and the proposed vintage-aware
design with lineage-driven incremental recomputation. On 1,435 cells compared
across two official channels, 33 differ; across two publication releases, 4,608
cell keys yield eight value differences in three domains. Of 123 classified
discrepancy events, 96 (78.05 per cent) are confirmed to arise from vintage,
methodology or granularity. On the frozen evidence workload the two treatments
that keep history recall all 38 published figures exactly, while overwrite and
single-source selection recall 14 of 38; the same result holds after a catalogue
restart and after a physically applied injected revision. Physical measurement
shows that at this size one revision costs one Iceberg commit rather than its
data: a single-value revision under overwrite changes 161 bytes of data but adds
23,915 bytes of metadata. The break-even point between incremental and full
recomputation is the subject of a frozen sweep that has not yet been executed.
\end{abstract}

\begin{highlights}
\item Differences between official SDG channels are characterised and attributed
to vintage, methodology or granularity by rules fixed in advance.
\item Indicator vintage is stored as an explicit dimension on an open table
format, keyed independently of table snapshots.
\item Four treatments are executed on an identical revision workload, logically
and physically, on the same hardware.
\item Recall of published figures, storage footprint and per-revision cost are
measured rather than asserted.
\item At indicator scale the cost of a revision is dominated by table metadata,
not by the revised value.
\end{highlights}

\begin{keywords}
official statistics \sep data vintage \sep reproducibility \sep data lineage
\sep lakehouse \sep SDG indicators
\end{keywords}

\maketitle

\section{Introduction}\label{sec:intro}

Indonesian SDG indicators are released by Badan Pusat Statistik (BPS) through
channels that run side by side. A web API serves dynamic table values that
follow database updates, while the annual SDG indicator compilation is published
as a document whose content is fixed on its release date
\citep{bps_tpb2024,bps_tpb2025}. Both are official, both are produced by the
same agency, and both are cited by the same users.

The two channels do not always agree. Comparing them cell by cell over the
snapshot taken for this study, 33 of 1,435 overlapping cells carry different
values. Comparing two consecutive editions of the published compilation, 4,608
aligned cell keys yield eight value differences spread over three domains,
namely energy, economy and ecology. Neither observation can be resolved by
declaring one channel wrong, because both are official products of the same
producer.

The causes can be separated. Applying eight classification rules fixed before
the data were inspected, 96 of 123 discrepancy events, or 78.05 per cent, are
confirmed to arise from one of three causes: a release vintage, a methodological
change, or a difference in granularity. The confirmed breakdown is 82
granularity events, 13 vintage events and one methodology event, while a further
19 events are rounding differences that the rules deliberately exclude from the
count. The remaining events stop at candidate status because the historical
snapshot needed to confirm them no longer exists.

That last sentence is the problem this paper addresses. Older values in the web
API are overwritten by newer ones, so a figure that was once published can no
longer be produced from its source. Revision policy for official statistics asks
for the opposite, namely the ability to show how an estimate changed from one
release to the next \citep{oecd_eurostat_revisions}. As long as that history
survives only inside a yearly document, the capability rests on a document
archive rather than on the system that produces the numbers.

Mechanisms for keeping history exist one layer down. Open table formats provide
snapshots and time travel \citep{iceberg_spec}, and incremental view maintenance
is a mature capability \citep{fabric_mlv}. The difficulty is what those
mechanisms record. A table snapshot records when a row was written, not when the
producer released the value. While data only accumulates the two time axes run
in parallel; when a source is revised retrospectively they separate, and the
table snapshot can no longer answer a question about the value as of a given
release.

This study therefore asks how much it costs to maintain SDG indicators that stay
reproducible and traceable when several official sources disagree, and at what
revision size incremental recomputation stops being cheaper than full
recomputation. Four sub-questions follow. \textbf{RQ1} asks which forms of
disagreement actually occur between official BPS sources across five domains,
and what share comes from vintage revision rather than methodology or
granularity. \textbf{RQ2} asks whether that disagreement can be represented as
indicator vintages on an open table format so that published values remain
addressable without storing a full copy of every release. \textbf{RQ3} asks what
time and space one revision costs, and where the incremental approach stops
paying. \textbf{RQ4} asks whether a previously reported figure can still be
produced exactly after its source has been revised, and under what conditions it
fails. RQ4 acts as an eligibility condition: a cheaper result that cannot
reproduce the old figure does not count as a gain.

The contributions are four. First, an empirical characterisation of disagreement
between official BPS sources across five domains, with classification rules
fixed in advance and measured cause proportions. Second, a storage design that
treats the indicator vintage as an explicit dimension of an open table format;
on the real workload, 38 observations from three releases project onto 14
indicator cells without losing provenance, and each per-release state is
reconstructed from the vintage key alone. Third, physical measurement of all
four treatments on identical hardware, including the finding that at this scale
one revision costs one table commit rather than its data. Fourth, a
reproducibility audit recomputed from the state of each treatment rather than
trusted from an earlier label.

The break-even point that RQ3 asks for is not reported here. The sweep that
locates it was frozen together with the rest of the experiment configuration and
is executed after this draft; Section~\ref{sec:status} states what remains.

\section{Related work and positioning}\label{sec:related}

\subsection{Computing SDG indicators from heterogeneous open data}

Computing SDG indicators from scattered open sources has been pursued for
several years. SustainGraph builds a knowledge graph that tracks progress and
interlinkage among SDG targets \citep{fotopoulou2022sustaingraph}, while
\citet{jiang2023discovery} propose a rule-based approach for discovering data
relevant to a given indicator. Both address discovery and linking, that is, how
a dataset is recognised as a candidate input for an indicator.

The closest work is SDG-KG \citep{benjira2025sdgkg}. It joins a metadata graph
derived from data sources with an SDG graph derived from the UN indicator
metadata repository, uses an LLM to assist schema alignment, and emits an
execution plan translated into Python code. When sources disagree, the conflict
is resolved by a trust score that combines spatial, temporal and contextual
similarity with source reliability and completeness, after which one source is
selected. The journal version extends the automated mapping between indicators
and open data \citep{benjira2025dke}.

Two limits apply to this line of work, and the authors state the first
themselves: SDG-KG operates at the metadata level and does not process data
flows, so it contains no performance evaluation, no scalability study and no
discussion of physical storage. The second limit is conceptual. Selecting one
source discards the values of the sources not selected. For open data of uneven
quality that is a reasonable decision. For official statistics it discards
precisely the information under study.

That consequence is measurable, and this study measures it. Treatment B2
replicates single-source selection using a frozen trust score. Ten of 14 cells
end up serving a value that does not come from the most recent vintage, and
under the injected revision workload only 8 of 28 revisions reach the serving
state, a propagation rate of 0.2857. This is not an implementation defect but a
property of the class of policies that selects one source before computation.

\subsection{Truth discovery and data fusion}

Truth discovery provides the formal basis for such a decision.
\citet{dong2015fusion} formulate data fusion as conflict resolution over many
sources, and subsequent surveys map the families of methods that estimate source
reliability while inferring the correct value
\citep{li2016truthsurvey,tbd2025truthsurvey}. The underlying assumption is
consistent across that literature: one hidden true value exists and the sources
are not all trustworthy.

The assumption does not hold for official statistics. When the BPS web API and
the published SDG compilation report different values for the same indicator,
region and year, no source is wrong. Both are issued by the same producer
through equally legitimate channels, and the difference originates in the
release vintage, a methodological change, or a difference in granularity. The
empirical evidence in this study supports that reading: 96 of 123 discrepancy
events are confirmed to arise from one of those three causes.

Truth discovery is therefore used as a comparator rather than as a foundation.
No novelty is claimed for its algorithms. Only the pattern is taken, so that it
can be run as treatment B2 and its consequences measured on the same workload as
the other treatments.

\subsection{Vintages and revisions in official statistics}

On the official statistics side the revision problem was formulated long ago.
The OECD and Eurostat guidelines set out how revisions are announced and
analysed \citep{oecd_eurostat_revisions}, and Statistics Canada publishes
real-time data tables that carry every revision of a data point over time
\citep{statcan_realtime}. Econometric work uses vintages to test forecast
stability \citep{ecb2007vintages}, and tooling continues to be developed for
analysing real-time series vintages.

Two terms from this literature are used directly here. A \emph{vintage} is the
state of a figure at the moment it was released, and a \emph{revision triangle}
is the matrix showing how the estimate for one period changes across subsequent
releases.

The limit lies in the object of study. This literature works on aggregates that
have already been published, usually small time series, in a statistical and
econometric setting. The pipeline that produces those figures is not the object,
so the question of how much time and space it takes to maintain their history
never arises. That question is the subject of this paper.

\subsection{Lakehouse versioning and incremental maintenance}

Mechanisms for storing and maintaining history are mature at the lakehouse
layer. Open table formats provide snapshots and time travel
\citep{iceberg_spec}, and systems research has compared their storage behaviour
experimentally \citep{jain2023lhbench,is2025lakehouse}. Incremental view
maintenance is likewise available as a product capability, with refreshes that
may be no-op, incremental or full \citep{fabric_mlv}.

The issue is not the availability of a mechanism but what the mechanism records,
as set out in Section~\ref{sec:intro}. The distinction is not merely conceptual.
In this study an Iceberg REST catalogue that kept its table registrations in
memory lost all of them when the node was restarted, although the data files
remained intact in object storage. The ability of the full-snapshot treatment to
recall old figures depends on snapshot metadata held in the catalogue, so that
ability disappeared without leaving any trace in the data. The proposed
treatment was unaffected at that point, because its recall rests on a
\texttt{(cell\_id, vintage\_id)} key stored as rows rather than on table
snapshots.

As with truth discovery, incremental view maintenance is used as a mechanism and
a comparator. No novelty is claimed here either.

\subsection{Provenance and pipeline reproducibility}

Reproducing a published figure depends on complete provenance. Transparent
provenance capture for data science pipelines has been addressed since URSPRUNG
\citep{rupprecht2020ursprung}, and the underlying principle was formulated
earlier as a requirement for computational research \citep{peng2011reproducible}.
This study follows that line and applies it to a specific chain, from source
manifest through vintage, transformation version, indicator cell and metric to
the table or figure. The completeness of that chain is measured rather than
assumed: 152 complete paths are formed across the four treatments, at a
completeness of 1.0000.

\subsection{Positioning}

Two mature literatures have not been brought together. Official statistics has a
settled notion of vintage but applies it to published aggregates without
accounting for computational cost. Data systems has complete machinery for
versioning and incremental maintenance but models disagreement between sources
as error to be resolved, a model that is wrong for official statistics.

This work sits between them, and differs from the closest prior work in three
ways. It operates at the physical layer, so what is measured is time, space and
the ability to reproduce a figure, rather than the ability to align schemas.
Differences between sources are treated as vintages to be preserved rather than
as conflicts from which one value must be chosen. Testing is carried out on
Indonesian national official statistics with revisions that actually occurred.

The claim boundary is stated in advance. No novelty is claimed for truth
discovery algorithms or for incremental view maintenance, and no claim is made
about scalability or performance, since all measurements run on a single small
node. What is claimed is the treatment of the indicator vintage as a first-class
object inside a lakehouse, together with the measurement of what maintaining it
costs.

\section{Study design}\label{sec:design}

\subsection{Treatments}

Four storage treatments are compared on identical workloads. Table~\ref{tbl:trt}
lists the state each keeps and the statements it issues to apply one revision;
the maintenance column matters because, as Section~\ref{sec:results} shows, the
number of statements dominates the measured time at this scale.

\begin{table}[t]
\caption{The four treatments. B3 is the proposal; B0--B2 are comparators.}\label{tbl:trt}
\begin{tabular*}{\tblwidth}{@{}p{0.06\textwidth}p{0.17\textwidth}p{0.18\textwidth}@{}}
\toprule
 & State kept & Write and maintenance per revision \\
\midrule
B0 & exactly one current row per cell & \texttt{MERGE}; expire snapshots to one \\
B1 & one full-state snapshot per release & \texttt{INSERT OVERWRITE} of the whole state; no expiry \\
B2 & single-source selection under a frozen trust score & \texttt{INSERT OVERWRITE} of the recomputed selection; expire snapshots to one \\
B3 & append-only store keyed by \texttt{(cell\_id, vintage\_id)} plus a materialised serving table & \texttt{INSERT} into the store, then \texttt{MERGE} of the dirty cells only; expire snapshots to one on both tables \\
\bottomrule
\end{tabular*}
\end{table}

The dirty set of B3 is derived from lineage edges, not from comparing values.
The trust score of B2 uses five weighted dimensions frozen before execution, and
observed values and conflict outcomes are barred from entering the score, so the
selection cannot be tuned by its own results.

\subsection{Workload and freeze}

The evidence workload consists of three release vintages, 38 observations and 14
indicator cells, and is the set of cells whose revisions are confirmed rather
than assumed. Revisions of a controlled size are injected on top of it, because
the size of a real revision cannot be chosen. Each injected revision changes one
value by exactly one unit of published precision, leaving every cell dimension
untouched; the mechanism is a dependency trigger, not a model of the empirical
revision distribution, so it is the number of revised cells that is controlled,
never the magnitude of the change.

The full experiment configuration was frozen before the main sweep: source
snapshots and retrieval dates, the indicator list per domain, the discrepancy
rules, the four treatments, the metric definitions, the injected revision sizes
and the resource limit. Each frozen value is recorded together with the checksum
of its evidence file at the moment of freezing.

All measurements run on a single node with 2 vCPU and 7.7~GiB of memory. Runs
are sequential and held by one operator, and every measurement point is repeated
three times and reported as the median. Timing may therefore be read only as a
comparison between treatments on identical hardware, never as a performance or
scalability claim.

\section{Results so far}\label{sec:results}

\subsection{Recall of published figures}

Table~\ref{tbl:res} reports the three measured outcomes per treatment. B1 and B3
recall all 38 published figures exactly; B0 and B2 recall 14 of 38, which is
what their contracts imply rather than a defect. The same outcome survives two
stress conditions applied physically: after the Iceberg catalogue was restarted,
and after an injected revision was written through each treatment's own write
path, where B1 and B3 recall 39 of 39 requests and B0 and B2 recall 14.

\begin{table}[t]
\caption{Measured outcomes per treatment on the frozen workload. Bytes are
medians over three repetitions; times are for a single-cell revision and include
the maintenance each treatment requires.}\label{tbl:res}
\begin{tabular*}{\tblwidth}{@{}lrrrr@{}}
\toprule
 & B0 & B1 & B2 & B3 \\
\midrule
Recall rate & 0.3684 & 1.0000 & 0.3684 & 1.0000 \\
Recall after restart & 14/38 & 38/38 & 14/38 & 38/38 \\
Recall after revision & 14/39 & 39/39 & 14/39 & 39/39 \\
Footprint (bytes) & 95,898 & 190,261 & 64,077 & 240,884 \\
Write time (s) & 9.818 & 5.442 & 5.308 & 13.388 \\
Total time (s) & 23.448 & 5.442 & 21.253 & 34.877 \\
Added bytes & 24,076 & 68,447 & 24,564 & 60,156 \\
\bottomrule
\end{tabular*}
\end{table}

\subsection{Space}

On the evidence workload B3 has the largest footprint, at 240,884 bytes against
190,261 for B1, because its serving table is materialised. Decomposing it shows
where the cost sits: the vintage store alone occupies 146,352 bytes, which is 23
per cent smaller than B1 while preserving the same recall. The cost of B3 on
this fixture therefore comes from serving materialisation and not from keeping
vintages. Iceberg metadata is comparable to the data it describes and exceeds it
for B0 and B3.

\subsection{Cost of one revision}

The timing results invert the expectation that the incremental treatment is
cheaper. On a single-cell revision B1 is the cheapest at 5.442~s, because it
issues one statement and requires no maintenance, while B3 is the most expensive
at 34.877~s across two statements plus snapshot expiry. At a state of 14 cells
the fixed cost of a statement exceeds its data work, so the ordering reflects
the number of statements rather than the number of rows touched.

The same effect appears in bytes. A single-value revision under B0 changes 161
bytes of data but adds 23,915 bytes of metadata, and B1 adds 68,447 bytes for
both a one-cell and a two-cell revision because it always rewrites the entire
state. B2 pays the full write cost even when it withholds the revision: on the
single-cell scenario it rewrites its selection table at a cost of 5.308~s and
24,564 bytes and still serves no new value, because the revision touches a cell
whose source loses on score. Selecting a single source therefore saves no write
cost; only information is lost.

Snapshot expiry is a real cost of the treatments that discard history. B0, B2 and
B3 spend between 14 and 23 seconds pruning snapshots on every revision, whereas
B1 needs none. Counting logical rows, as an earlier stage of this work did,
makes that cost invisible.

\section{Current status and limitations}\label{sec:status}

The results above are measured on a fixture of 38 observations across 14 cells,
with revisions of one and two cells. Neither the time ordering nor the byte
ordering may be extrapolated from that size. In particular, the break-even point
that RQ3 asks for cannot be observed on this fixture, precisely because the
fixed cost per statement dominates the data work; that is an empirical finding
about the measurement design, not an answer to RQ3.

The sweep that answers RQ3 is frozen and pending. It runs on a province panel of
5,378 indicator cells, roughly 384 times the base state of the fixture, with
revision sizes of 1, 2, 4, 8, 16 and 38 cells, where the largest point is one
full year for every province of a single indicator served by a single source. At
that base-state size the full rewrite of B1 grows with the size of the state
while the work of B3 grows with the number of revised cells, which is the
difference that decides whether a break-even point exists. Its result, together
with the analysis of the case in which an old figure cannot be reproduced, will
complete the answer to RQ3 and RQ4.

Two further limits are stated in advance. Data volume in this study is small,
which follows from the domain rather than from the design: SDG indicators at
national and province granularity are never large, and the questions asked
concern correctness, traceability and maintenance cost rather than throughput.
Second, the results are specific to one national statistical office and to five
SDG domains, and are not generalised to all 17 goals or to other countries.

\section*{Data and code availability}

The source manifests, checksums, classification rules, frozen experiment
configuration, treatment implementations and every processed result table are
kept in the project repository. Raw source payloads are not committed;
every figure in this paper is reproducible from the committed manifests without
re-downloading the sources.

\printcredits

\bibliographystyle{cas-model2-names}
\bibliography{refs}

\end{document}
